\chapter{Deliverable}
\label{chap:deliverable}

\section{Software Solution}
\label{section:software-solution}

The delivered SpotOn system comprises a fully functional web-based tracking dashboard and a robust computer vision backend pipeline. The software deliverables successfully unify real-time AI inference with an accessible user interface tailored for security and facility management. The solution is divided into three core pillars:

\subsection{Frontend Web Application}
Developed using React with Vite and utilizing Tailwind CSS, the frontend provides an intuitive, high-performance command center. The key pages delivered are detailed below.
\newpage

\subsubsection{Landing Page}
\begin{figure}[H]
    \centering
    \includegraphics[width=\textwidth]{jubjones/UI/LANDING.jpeg}
    \caption{SpotOn Landing Page Interface}
    \label{fig:deliv-landing}
\end{figure}
The Landing Page is the initial entry point of the application. It provides users with a high-level overview of the SpotOn system's capabilities and serves as a gateway to securely access the core tracking environments.
\newpage

\subsubsection{Live Video Feed}
\begin{figure}[H]
    \centering
    \includegraphics[width=\textwidth]{jubjones/UI/VIDEO_FEED.jpeg}
    \caption{Live Video Feed Interface with Real-Time Data}
    \label{fig:deliv-videofeed}
\end{figure}
The Live Video Feed serves as the core operational dashboard. It simultaneously renders real-time tracking streams adjacent to a 2D topographical floor plan. The interface actively translates individual bounding boxes from the video perspective into persistent, unified coordinates directly on the map.
\newpage

\subsubsection{Analytics Page}
\begin{figure}[H]
    \centering
    \includegraphics[width=\textwidth]{jubjones/UI/ANALYTIC_01.png}
    \caption{Analytics Dashboard showing insights and trends}
    \label{fig:deliv-analytics}
\end{figure}
The Analytics Page is a dedicated reporting module summarizing historical operational metrics. It visualizes critical data such as zone traffic densities, peak activity times, localized dwell durations, and trajectory heatmaps based on the securely stored tracking data.
\newpage

\subsubsection{Help Page}
\begin{figure}[H]
    \centering
    \includegraphics[width=\textwidth]{jubjones/UI/HELP.png}
    \caption{Help Page providing documentation and system guides}
    \label{fig:deliv-help}
\end{figure}
The Help Page acts as a comprehensive support module containing user guides, system documentation, and troubleshooting steps to assist operators in navigating the dashboard and effectively understanding the AI tracking outputs.
\newpage

\subsection{Real-Time Backend Services}
A FastAPI-driven service architecture operates as the central controller, designed for high throughput and low latency.
\begin{itemize}
    \item \textbf{Data Ingestion \& Streaming:} Efficient handling of continuous video feeds connected through internal deployment channels, parsing streams frame-by-frame.
    \item \textbf{Asynchronous Dispatch:} Using WebSockets to instantaneously broadcast bounding box coordinates, unique tracking IDs, and system status markers back to the connected frontend clients without blocking the main event loops.
    \item \textbf{State Persistence:} Integrating PostgreSQL and TimescaleDB to securely log tracking metadata and system events, ensuring the analytics pipeline retains temporally accurate data.
\end{itemize}

\vspace{\fill}
\newpage

\begin{figure}[H]
    \centering
    \includegraphics[width=\textwidth]{jubjones/backend_api.png}
    \caption{Interactive FastAPI Swagger documentation detailing system endpoints.}
    \label{fig:deliv-backend-api}
\end{figure}

\subsection{AI Models \& Inference Pipeline}
The intelligence driving the physical tracking relies on a multi-stage computer vision workflow integrated directly into the backend.
\begin{itemize}
    \item \textbf{Object Detection \& Spatial Tracking:} Utilizing YOLO-based models for high-speed pedestrian detection, tightly coupled with SORT-driven algorithms to maintain short-term, single-camera spatial continuity.
    \item \textbf{Cross-Camera Re-Identification (Re-ID):} An implementation of a Torchreid-based neural network architecture capable of extracting deep appearance embeddings using an OSNet-AIN backbone. This enables the system to re-identify subjects across disjointed camera views by comparing feature vectors, mitigating the limitations of blind spots.
    \item \textbf{Experiment Tracking \& Model Management:} Throughout the development lifecycle, MLflow was utilized to systematically log all model training experiments, track hyperparameters, and compare evaluation metrics across different iterations of the AI models. This ensured rigorous engineering practices and reproducibility before finalizing the weights for deployment.
\end{itemize}

\vspace{\fill}
\newpage

\begin{figure}[H]
    \centering
    \includegraphics[width=\textwidth]{jubjones/mlflow_experiment.png}
    \caption{MLflow dashboard displaying experiment tracking and model evaluation metrics.}
    \label{fig:deliv-mlflow}
\end{figure}

\subsection{Continuous Integration and Deployment (CI/CD)}
To ensure code quality, prevent regressions, and streamline the development lifecycle, SpotOn employs automated Continuous Integration (CI) pipelines using GitHub Actions for both the frontend and backend repositories. These pipelines are triggered automatically upon every pull request and push to the main branches.

\begin{itemize}
    \item \textbf{Backend CI Pipeline:} The backend workflow is optimized for speed and isolation. It utilizes \texttt{uv}, an extremely fast Python package and project manager, to provision Python 3.12 environments and install dependencies. To keep the cloud runners efficient, the pipeline specifically requests CPU-only wheels for heavy machine learning libraries (like PyTorch) and avoids installing GPU-specific CUDA toolkits. Finally, it executes an extensive suite of over 160 isolated unit and integration tests using \texttt{pytest}, ensuring API contracts, Re-ID logic, and database operations remain fully functional.
    \item \textbf{Frontend CI Pipeline:} The frontend workflow utilizes Node.js and \texttt{npm} to ensure UI stability. The pipeline executes a strict progression of quality checks: running static code linting via ESLint to maintain styling conventions, enforcing rigorous TypeScript type-checking to catch compilation errors, and ultimately executing the suite of unit and component tests to validate UI rendering and state management logic.
\end{itemize}

\vspace{\fill}
\newpage

\begin{figure}[H]
    \centering
    \includegraphics[width=\textwidth]{jubjones/backend_ci.png}
    \caption{GitHub Actions dashboard illustrating successful automated CI/CD pipeline executions.}
    \label{fig:deliv-github-actions}
\end{figure}

The reliability and performance of the SpotOn system were validated through a series of rigorous tests covering both frontend interface interactions and backend service logic. The following tables summarize the test cases and their respective execution statuses.

\subsection{Frontend Component Testing}
The frontend test suite focuses on page rendering, user navigation, real-time data visualization, and accessibility across various device viewports.

\begin{longtable}{|c|p{4cm}|p{4cm}|p{4cm}|c|}
\caption{SpotOn Frontend Test Summary} \label{tab:frontend_tests} \\
\hline
\textbf{\#} & \textbf{Test Case Name} & \textbf{Description} & \textbf{Expected Result} & \textbf{Status} \\ \hline
\endfirsthead
\multicolumn{5}{c}%
{{\bfseries \tablename\ \thetable{} -- continued from previous page}} \\
\hline
\textbf{\#} & \textbf{Test Case Name} & \textbf{Description} & \textbf{Expected Result} & \textbf{Status} \\ \hline
\endhead
\hline \multicolumn{5}{|r|}{{Continued on next page}} \\ \hline
\endfoot
\hline
\endlastfoot
1 & Landing Page Rendering & Visit the landing page & Displays branding, status cards, and action buttons & Passed \\ \hline
2 & Header Navigation & Click header links (Help, About, Analytics) & Navigates to correct pages and returns via logo & Passed \\ \hline
3 & Backend Status & Check connection indicator & Shows ``Connected'' with green indicator when backend is healthy & Passed \\ \hline
4 & Mobile Navigation & Set viewport to mobile size & Displays mobile navigation and adapted layout & Passed \\ \hline
5 & Environment List & Visit /environments page & Displays all pre-configured monitoring environments & Passed \\ \hline
6 & Environment Selection & Click an environment card & Environment is highlighted and ``Continue'' button appears & Passed \\ \hline
7 & Environment Details & Click environment for details & Shows camera list and specific configuration for that environment & Passed \\ \hline
8 & Navigation to Tracking & Select environment and click Continue & Navigates to /group-view with correct environment ID & Passed \\ \hline
9 & Multi-Camera Grid & Load tracking view for environment & All associated cameras are displayed in a 2x2 grid layout & Passed \\ \hline
10 & Camera Focus Mode & Click on a specific camera & Camera expands to focus mode, others are minimized & Passed \\ \hline
11 & Real-time Data Feed & Dispatch tracking-update event & Bounding boxes are rendered on relevant camera feeds & Passed \\ \hline
12 & Layout Controls & Change grid layout to 1x4 & Grid adjusts to single column layout with 4 rows & Passed \\ \hline
13 & Playback Controls & Click pause/play buttons & UI state reflects playback status (paused/playing) & Passed \\ \hline
14 & WebSocket Recovery & Simulate WebSocket disconnect & Shows ``Reconnecting'' status and resumes on reconnect & Passed \\ \hline
15 & Camera Error Handling & Mock 404 for camera frame & Specific camera shows ``Frame not available'' error state & Passed \\ \hline
16 & Person Highlighting & Click on a person bounding box & Bounding box is highlighted across all cameras and map & Passed \\ \hline
17 & Person Details Panel & Select a person & Details panel opens with ID, confidence, and timeline & Passed \\ \hline
18 & Cross-Camera Tracking & Move person from c01 to c02 & Person disappears from c01 and appears in c02 with same ID & Passed \\ \hline
19 & High-Frequency Load & Send 100 rapid tracking updates & UI remains responsive and renders tracks without lag & Passed \\ \hline
20 & Map Rendering & Load group view & Tracking map is visible with proper dimensions & Passed \\ \hline
21 & Camera Markers & Check map content & Camera positions are correctly marked on the map & Passed \\ \hline
22 & Person Markers & Update tracking data & Person marker appears on map at correct coordinates & Passed \\ \hline
23 & Map Interaction & Zoom and pan the map & Map responds to mouse wheel and drag interactions & Passed \\ \hline
24 & Analytics Rendering & Visit /analytics page & Displays total detections, active tracks, and uptime metrics & Passed \\ \hline
25 & Chart Visualization & Check analytics content & Time-series charts are rendered and visible & Passed \\ \hline
26 & Data Export & Click export button & Displays format options (CSV, JSON, PDF) & Passed \\ \hline
27 & Settings Page & Visit /settings page & Shows camera, detection, alert, and system sections & Passed \\ \hline
28 & Threshold Adjustment & Fill confidence threshold input & Value is updated and persistence logic is triggered & Passed \\ \hline
29 & Camera Toggling & Toggle a camera switch & Camera state updates and reflects in UI & Passed \\ \hline
30 & Search Functionality & Enter query in search bar & Filters help topics based on search text & Passed \\ \hline
31 & FAQ Interaction & Click an FAQ question & Question expands to show answer and collapses on second click & Passed \\ \hline
32 & Section Navigation & Click section links & Navigates to Getting Started, Features, Troubleshooting, etc. & Passed \\ \hline
33 & API Status Check & Call getSystemHealth() & Returns healthy status from backend & Passed \\ \hline
34 & Task Management & Call startProcessingTask() & Task is created successfully with WebSocket URL & Passed \\ \hline
35 & WebSocket Service & Connect to WebSocket & Connection established and message handlers registered & Passed \\ \hline
36 & Error Boundary & Trigger component crash & Renders fallback error UI instead of breaking the app & Passed \\ \hline
37 & Accessibility & Check ARIA labels and roles & Interactive elements have accessible names and semantic roles & Passed \\ \hline
38 & Responsive Utils & Change screen dimensions & Breakpoints correctly detect mobile/tablet/desktop states & Passed \\ \hline
\end{longtable}

\subsection{Backend Service Testing}
The backend test report covers API availability, environment management, real-time detection tasks, and critical AI components such as Re-Identification and Identity merging.

\begin{longtable}{|c|p{1.5cm}|p{3.5cm}|p{3.5cm}|p{3.5cm}|c|}
\caption{SpotOn Backend Test Summary} \label{tab:backend_tests} \\
\hline
\textbf{\#} & \textbf{ID} & \textbf{Name} & \textbf{Description} & \textbf{Expected Result} & \textbf{Status} \\ \hline
\endfirsthead
\multicolumn{6}{c}%
{{\bfseries \tablename\ \thetable{} -- continued from previous page}} \\
\hline
\textbf{\#} & \textbf{ID} & \textbf{Name} & \textbf{Description} & \textbf{Expected Result} & \textbf{Status} \\ \hline
\endhead
\hline \multicolumn{6}{|r|}{{Continued on next page}} \\ \hline
\endfoot
\hline
\endlastfoot
1 & HLT-001 & Root Welcome & GET \texttt{/} & Returns 200 and welcome message & Passed \\ \hline
2 & HLT-002 & Basic Health & GET \texttt{/health} & Returns 200 with status field & Passed \\ \hline
3 & HLT-003 & Deep Health & GET \texttt{/health/deep} & Returns 200/503/500 with component status & Passed \\ \hline
4 & HLT-004 & Liveness Probe & GET \texttt{/health/live} & Returns 200 ``alive'' & Passed \\ \hline
5 & HLT-005 & WS Module & GET \texttt{/ws/health} & Returns 200 status & Passed \\ \hline
6 & ENV-001 & List environments & GET \url{/api/v1/environments/} & Returns 200 list & Passed \\ \hline
7 & ENV-002 & Env metadata & GET valid ID & Returns 200 details & Passed \\ \hline
8 & ENV-003 & Unknown Env & GET invalid ID & Returns 404 & Passed \\ \hline
9 & ENV-004 & List cameras & GET cameras & Returns 200 camera list & Passed \\ \hline
10 & ENV-005 & List zones & GET zones & Returns 200 zone list & Passed \\ \hline
11 & ENV-006 & Create session & POST no data & Returns 400 when no data exists & Passed \\ \hline
12 & ENV-007a & Get pref & GET preferences & Returns 200 with user\_id & Passed \\ \hline
13 & ENV-007b & Update pref & PUT preferences & Updates user data successfully & Passed \\ \hline
14 & TSK-001 & Create task & POST valid body & Returns 202 accepted & Passed \\ \hline
15 & TSK-002 & Task (Invalid) & POST empty body & Returns 422 validation error & Passed \\ \hline
16 & TSK-003 & Task Details & GET unknown ID & Returns 404 not found & Passed \\ \hline
17 & TSK-004 & Poll status & GET unknown ID & Returns 404 for unknown task & Passed \\ \hline
18 & TSK-005 & List tasks & GET all tasks & Returns 200 list of tasks & Passed \\ \hline
19 & TSK-006 & Stop task & DELETE unknown & Returns 404 for unknown task & Passed \\ \hline
20 & ANL-001 & Dashboard sum & GET dashboard & Dashboard responds successfully & Passed \\ \hline
21 & ANL-002 & RT metrics & GET real-time & Metrics respond successfully & Passed \\ \hline
22 & ANL-003 & Behavior anal & POST analysis & Analysis responds successfully & Passed \\ \hline
23 & ANL-004 & Hist summary & POST summary & Summary responds successfully & Passed \\ \hline
24 & ANL-005 & System stats & GET statistics & Statistics respond successfully & Passed \\ \hline
25 & PLY-001 & Pause proc & POST pause & Exercises routing and pauses task & Passed \\ \hline
26 & PLY-002 & Resume proc & POST resume & Exercises routing and resumes task & Passed \\ \hline
27 & PLY-003 & Seek in stream & POST seek & Exercises routing and seeks offset & Passed \\ \hline
28 & STR-001 & MJPEG Feed & GET multipart & Returns valid multipart JPEG stream & Passed \\ \hline
29 & STR-002 & Disconnection & Stream disconnect & Server closes connection cleanly & Passed \\ \hline
30 & STR-003 & Invalid Task & Request stream & Handles invalid task ID correctly & Passed \\ \hline
31 & WS-001 & Telemetry & WS tracking & WebSocket connection accepted & Passed \\ \hline
32 & WS-002 & Frame Meta & WS frames & WebSocket connection accepted & Passed \\ \hline
33 & WS-003 & System Events & WS system & WebSocket connection accepted & Passed \\ \hline
34 & WS-004 & Live Analytics & WS analytics & WebSocket connection accepted & Passed \\ \hline
35 & WS-005 & Drop handling & WS disconnect & Abrupt disconnect handled without crash & Passed \\ \hline
36 & WS-006 & Client message & Send JSON focus & JSON processed correctly by server & Passed \\ \hline
37 & PERF-API-001 & API Latency & Measure p99 & Latencies within acceptable thresholds & Passed \\ \hline
38 & PERF-API-002 & Throughput & Baseline RPS & Meets required sequential throughput & Passed \\ \hline
39 & PERF-API-003 & Error Rate & Normal load & 4xx/5xx fraction remains minimal & Passed \\ \hline
40 & PERF-RT-001 & Proc Latency & Frame handler & Overhead remains under 30ms & Passed \\ \hline
41 & PERF-RT-002 & Init Time & MJPEG TTFB & Stream starts in under 500ms & Passed \\ \hline
42 & PERF-RT-003 & WS Latency & Round-trip & Latency remains under target threshold & Passed \\ \hline
43 & PERF-RT-004 & Drop Rate & Consumption & Minimal frame drops under load & Passed \\ \hline
44 & PERF-RT-000 & Glass-to-Glass & End-to-end & Orchestration within design spec & Passed \\ \hline
45 & REID-001 & Feature Ext & Extraction & Produces 512-dim embeddings & Passed \\ \hline
46 & REID-002 & Handoff Match & Matching & Match found across camera views & Passed \\ \hline
47 & REID-003 & Handoff Expiry & Expiry logic & Records expire after time window & Passed \\ \hline
48 & REG-001 & Assignment & Identity alloc & Assigns global ID to local track & Passed \\ \hline
49 & REG-002 & Identity Merge & ID Merging & Merging global IDs updates all tracks & Passed \\ \hline
50 & REG-003 & Shared Check & Identity check & Detects subjects across multiple cams & Passed \\ \hline
\end{longtable}

\section{System Performance Benchmarks}
\label{section:system-performance-benchmarks}

To validate the theoretical capabilities outlined in the system architecture and AI design chapters, the SpotOn system underwent rigorous empirical benchmarking. The following results quantify both the accuracy of the underlying AI models and the real-time processing capabilities of the backend architecture. 

\subsection{AI Accuracy Metrics}
The primary objective of the AI pipeline is to maintain high accuracy in dynamic, crowded scenarios. The detection model (YOLOv26m) and Re-Identification model (Torchreid/OSNet-AIN) were evaluated against the MTMMC validation dataset. The key metrics reflect the system's ability to minimize identity fragmentation and accurately localize subjects.



\begin{table}[H]
    \centering
    \caption{Detailed Tracking Performance per Camera (TensorRT Optimized, frame-skip = 2)}
    \label{tab:detailed_tracking}
    \resizebox{\columnwidth}{!}{%
    \begin{tabular}{|l|c|c|c|c|c|c|c|c|c|}
    \hline
    \textbf{Camera} & \textbf{mAP@50} & \textbf{Recall} & \textbf{Precision} & \textbf{MOTA} & \textbf{MOTP} & \textbf{IDF1} & \textbf{ID Sw.} & \textbf{Re-ID mAP} & \textbf{Rank-1} \\
    \hline
    \multicolumn{10}{|c|}{\textbf{Factory Environment}} \\
    \hline
    c09 & 74.5\% & 82.2\% & 82.7\% & 63.9\% & 9.4\% & 68.6\% & 6 & 71.6\% & 78.3\% \\
    c12 & 63.4\% & 74.8\% & 75.6\% & 52.4\% & 15.5\% & 62.8\% & 1 & 60.1\% & 66.8\% \\
    c13 & 67.2\% & 78.6\% & 78.3\% & 57.4\% & 13.8\% & 62.8\% & 1 & 65.1\% & 71.8\% \\
    c16 & 68.6\% & 77.3\% & 82.3\% & 59.7\% & 11.5\% & 58.9\% & 42 & 67.4\% & 74.1\% \\
    \hline
    \textbf{Factory Avg.} & \textbf{68.4\%} & \textbf{78.2\%} & \textbf{79.7\%} & \textbf{58.4\%} & \textbf{12.5\%} & \textbf{63.3\%} & \textbf{50} & \textbf{66.1\%} & \textbf{72.8\%} \\
    \hline
    \multicolumn{10}{|c|}{\textbf{Campus Environment}} \\
    \hline
    c01 & 70.3\% & 82.1\% & 75.7\% & 57.4\% & 19.3\% & 59.0\% & 48 & 65.1\% & 71.8\% \\
    c02 & 56.4\% & 71.7\% & 72.2\% & 47.2\% & 18.7\% & 60.7\% & 3 & 54.9\% & 61.6\% \\
    c03 & 50.7\% & 69.1\% & 66.5\% & 39.8\% & 23.4\% & 55.3\% & 15 & 47.5\% & 54.2\% \\
    c05 & 65.0\% & 77.1\% & 78.0\% & 56.1\% & 17.0\% & 65.7\% & 0 & 63.8\% & 70.5\% \\
    \hline
    \textbf{Campus Avg.} & \textbf{60.6\%} & \textbf{75.0\%} & \textbf{73.1\%} & \textbf{50.1\%} & \textbf{19.6\%} & \textbf{60.2\%} & \textbf{66} & \textbf{57.8\%} & \textbf{64.5\%} \\
    \hline
    \end{tabular}%
    }
\end{table}

\subsection{Operations and Real-Time Pipeline Metrics}
To guarantee situational awareness without disruptive lag, the system's streaming and inference loops were load-tested to measure latency, throughput, and consistency. These metrics determine the true ``real-time'' capability of the SpotOn dashboard.

\begin{table}[H]
    \centering
    \caption{Average Per-Frame Processing Latency (TensorRT GPU)}
    \label{tab:ops_performance_metrics}
    \begin{tabular}{|l|c|}
    \hline
    \textbf{Pipeline Stage} & \textbf{Avg. Latency (ms)} \\ \hline
    Object Detection (YOLOv26m TensorRT) & 11.9 \\ \hline
    Data Association (ByteTrack) & 1.2 \\ \hline
    Feature Extraction (OSNet) & 36.3 \\ \hline
    \textbf{Total End-to-End Processing} & \textbf{49.4} \\ \hline
    \textbf{Estimated FPS Cap} & \textbf{20.3 FPS} \\ \hline
    \end{tabular}
\end{table}